\BourbakiExercisesSection

\begin{enumerate}
\def\labelenumi{\arabic{enumi}.}
\item
  在一个 n 阶矩阵 U 中，如果对于每个指标 i，将指标 i 的列替换为指标 \(\neq i\) 的列之和，则所得矩阵的行列式等于 \((-1)^{n-1}(n-1)\det U\) 。如果在 U 中，对于每个 i，从指标 i 的列中减去指标 \(\neq i\) 的列之和，则所得的行列式等于 \(-(n-2)2^{n-1}\det U\) 。
\item
  设 \(\Delta = \det (\alpha_{ij})\) 是一个 \(n\) 阶矩阵的行列式；对于
\end{enumerate}

\[
1 \leqslant i \leqslant n - 1
\]

和 \(1 \leqslant j \leqslant n - 1\) ，我们记作

\[
\beta_ {i j} = \left| \begin{array}{c c} \alpha_ {\mathbf {1} j} & \alpha_ {\mathbf {1}, j + \mathbf {1}} \\ \alpha_ {i + \mathbf {1}, j} & \alpha_ {i + \mathbf {1}, j + \mathbf {1}} \end{array} \right|.
\]

证明 n - 1 阶行列式 \(\det(\beta_{ij})\) 等于

\[
\alpha_ {1 2} \alpha_ {1 3} \dots \alpha_ {1, n - 1} \Delta .
\]

\begin{enumerate}
\def\labelenumi{\arabic{enumi}.}
\setcounter{enumi}{2}
\tightlist
\item
  证明恒等式
\end{enumerate}

\[
\left| \begin{array}{cccccc} x _ {1} & y _ {1} & \alpha_ {1 3} & \alpha_ {1 4} & \ldots & \alpha_ {1 n} \\ \lambda_ {1} x _ {2} & x _ {2} & y _ {2} & \alpha_ {2 4} & \ldots & \alpha_ {2 n} \\ \lambda_ {1} \lambda_ {2} x _ {3} & \lambda_ {2} x _ {3} & x _ {3} & y _ {3} & \ldots & \alpha_ {3 n} \\ \cdot & \cdot & \cdot & \cdot & \cdot & \cdot \\ \lambda_ {1} \lambda_ {2} \ldots \lambda_ {n - 1} x _ {n} & \lambda_ {2} \ldots \lambda_ {n - 1} x _ {n} & \lambda_ {3} \ldots \lambda_ {n - 1} x _ {n} & \lambda_ {4} \ldots \lambda_ {n - 1} x _ {n} & \ldots & x _ {n} \\ \end{array} \right| = (x _ {1} - \lambda_ {1} y _ {1}) (x _ {2} - \lambda_ {2} y _ {2}) \ldots (x _ {n - 1} - \lambda_ {n - 1} y _ {n - 1}) x _ {n}.
\]

推导出以下恒等式：

\[
\left| \begin{array}{c c c c c} 1 & 1 & 1 & \dots & 1 \\ b _ {1} & a _ {1} & a _ {1} & \dots & a _ {1} \\ b _ {1} & b _ {2} & a _ {2} & \dots & a _ {2} \\ \cdot & \cdot & \cdot & \cdot & \cdot \\ b _ {1} & b _ {2} & b _ {3} & \dots & a _ {n} \end{array} \right| = (a _ {1} - b _ {1}) (a _ {2} - b _ {2}) \dots (a _ {n} - b _ {n})
\]

\[
\left| \begin{array}{c c c c c} a _ {1} b _ {1} & a _ {1} b _ {2} & a _ {1} b _ {3} & \ldots & a _ {1} b _ {n} \\ a _ {1} b _ {2} & a _ {2} b _ {2} & a _ {2} b _ {3} & \ldots & a _ {2} b _ {n} \\ a _ {1} b _ {3} & a _ {2} b _ {3} & a _ {3} b _ {3} & \ldots & a _ {3} b _ {n} \\ \cdot \cdot & \cdot \cdot & \cdot \cdot & \cdot \cdot & \cdot \cdot \\ a _ {1} b _ {n} & a _ {2} b _ {n} & a _ {3} b _ {n} & \ldots & a _ {n} b _ {n} \end{array} \right| = a _ {1} b _ {n} (a _ {2} b _ {1} - a _ {1} b _ {2}) \ldots (a _ {n} b _ {n - 1} - a _ {n - 1} b _ {n})
\]

\[
\left| \begin{array}{cccccc} a _ {2} a _ {3} \ldots a _ {n} & a _ {3} a _ {4} \ldots a _ {n} b _ {1} & a _ {4} \ldots a _ {n} b _ {1} b _ {2} & \ldots & b _ {1} b _ {2} b _ {3} \ldots b _ {n - 1} \\ b _ {2} b _ {3} \ldots b _ {n} & a _ {3} a _ {4} \ldots a _ {n} a _ {1} & a _ {4} \ldots a _ {n} a _ {1} b _ {2} & \ldots & a _ {1} b _ {2} b _ {3} \ldots b _ {n - 1} \\ a _ {2} b _ {3} \ldots b _ {n} & b _ {3} b _ {4} \ldots b _ {n} b _ {1} & a _ {4} \ldots a _ {n} a _ {1} a _ {2} & \ldots & a _ {1} a _ {2} b _ {3} \ldots b _ {n - 1} \\ \cdot & \cdot & \cdot & \cdot & \cdot \\ a _ {2} a _ {3} \ldots b _ {n} & a _ {3} a _ {4} \ldots b _ {n} b _ {1} & a _ {4} \ldots b _ {n} b _ {1} b _ {2} & \ldots & a _ {1} a _ {2} a _ {3} \ldots a _ {n - 1} \\ \end{array} \right| = (a _ {1} a _ {2} \ldots a _ {n} - b _ {1} b _ {2} \ldots b _ {n}) ^ {n - 1}.
\]

\[
\left| \begin{array}{c c c c c c} x & a _ {1} & a _ {2} & \ldots & a _ {n - 1} & 1 \\ a _ {1} & x & a _ {2} & \ldots & a _ {n - 1} & 1 \\ a _ {1} & a _ {2} & x & \ldots & a _ {n - 1} & 1 \\ . & . & . & . & . & . \\ a _ {1} & a _ {2} & a _ {3} & \ldots & x & 1 \\ a _ {1} & a _ {2} & a _ {3} & \ldots & a _ {n} & 1 \end{array} \right| = (x - a _ {1}) (x - a _ {2}) \ldots (x - a _ {n})
\]

\[
\begin{array}{r l} & {\left| \begin{array}{l l l l l} x & a _ {1} & a _ {2} & \ldots & a _ {n} \\ a _ {1} & x & a _ {2} & \ldots & a _ {n} \\ a _ {1} & a _ {2} & x & \ldots & a _ {n} \\ \cdot & \cdot & \cdot & \cdot & \cdot \\ a _ {1} & a _ {2} & a _ {3} & \ldots & x \end{array} \right|} \\ & {\qquad = (x + a _ {1} + a _ {2} + \dots + a _ {n}) (x - a _ {1}) (x - a _ {2}) \ldots (x - a _ {n})} \end{array}
\]

（将最后一个行列式化简为前一个行列式）。

\begin{enumerate}
\def\labelenumi{\arabic{enumi}.}
\setcounter{enumi}{3}
\tightlist
\item
  计算行列式
\end{enumerate}

\[
\Delta_ {n} = \left| \begin{array}{c c c c c} a _ {1} + b _ {1} & b _ {1} & b _ {1} & \ldots & b _ {1} \\ b _ {2} & a _ {2} + b _ {2} & b _ {2} & \ldots & b _ {2} \\ \cdot & \cdot & \cdot & \cdot & \cdot \\ b _ {n} & b _ {n} & b _ {n} & \ldots & a _ {n} + b _ {n} \end{array} \right|
\]

（将 \(\Delta_{n}\) 用 \(\Delta_{n-1}\) 表出）。

\begin{enumerate}
\def\labelenumi{\arabic{enumi}.}
\setcounter{enumi}{4}
\tightlist
\item
  若 \(a_i\) 与 \(b_j\) 是交换域的元素，使得 \(a_i + b_j \neq 0\) （对每一对有序指标 \((i,j)\)），证明
\end{enumerate}

\[
\det \left(\frac {1}{a _ {i} + b _ {j}}\right) = \frac {\sum_ {i <   j} (a _ {j} - a _ {i}) (b _ {j} - b _ {i})}{\sum_ {i , j} (a _ {i} + b _ {j})}
\]

（“柯西行列式”）。

\begin{enumerate}
\def\labelenumi{\arabic{enumi}.}
\setcounter{enumi}{5}
\tightlist
\item
  证明：若 X 是 n 行 m 列的矩阵，Y 是 p 行 n 列的矩阵，且 Z = YX 是 p 行 m 列的乘积矩阵，则 Z 的 q 阶子式在 n \textless{} q 时为零；若 \(q \leqslant n\) ，它们由
\end{enumerate}

\[
\det (Z _ {L, H}) = \sum_ {K} \det (Y _ {L, K}) \det (X _ {K, H})
\]

其中 K 遍历 \([1, n]\) 的具有 q 个元素的子集所成的集合（利用 §7 第 2 号中的公式 (3)）。

\begin{enumerate}
\def\labelenumi{\arabic{enumi}.}
\setcounter{enumi}{6}
\item
  设 \(\Delta = \det(\alpha_{ij})\) 是 n 阶矩阵 U 的行列式；对每个指标 i（ \(1 \leqslant i \leqslant n\) ），令 \(\Delta_i\) 表示将 U 中的元素 \(\alpha_{ij}\) 乘以 \(\beta_j\) （对 \(1 \leqslant j \leqslant n\)）后所得的行列式。证明这 n 个行列式 \(\Delta_i\) ( \(1 \leqslant i \leqslant n\) ）之和等于 \((\beta_1 + \beta_2 + \cdots + \beta_n)\Delta\) （按指标 i 的行展开 \(\Delta_i\) ）。
\item
  设 \(\Delta = \det(\alpha_{ij})\) 是矩阵 \(U\) （\(n\) 阶）的行列式， \(\sigma\) 是属于 \(\mathfrak{S}_n\) 的一个置换；对每个指标 \(i\) ( \(1 \leqslant i \leqslant n\) )，令 \(\Delta_i\) 为把元素 \(\alpha_{ij}\) 在 \(U\) 中换成 \(\alpha_{i,\sigma(j)}\) （对 \(1 \leqslant j \leqslant n\)）所得的行列式；若 \(p\) 是在置换 \(\sigma\) 下不变的指标的个数，证明这 \(n\) 个行列式 \(\Delta_i\) ( \(1 \leqslant i \leqslant n\) ）之和等于 \(p\Delta\) （方法与习题 7 相同）。
\item
  设 A 为 n 阶方阵，B 为 A 的一个具有 p 行 q 列的子矩阵，C 为把 A 中属于 B 的每个元素乘以同一标量 \(\alpha\) 后所得的矩阵。证明 det C 的总展开式中的每一项都等于 det A 的总展开式中对应项乘以一个形如 \(\alpha^{r}\) 的标量，其中 \(r \geqslant p + q - n\) ，而 r 取决于所考虑的项（对 det C 作适当的拉普拉斯展开）。
\item
  设 \(\Gamma, \Delta\) 是两个矩阵 \(U\) , \(V\) （\(n\) 阶）的行列式；设 H 和 K 是 \([1, n]\) 的任意两个具有 \(p\) 个元素的子集，而 \((i_k)\) （相应地 \((j_k)\) ）是将 H（相应地，K）的元素按递增顺序排列所得的序列；令 \(\Gamma_{H,K}\) 为在 \(U\) 中把指标为 \(i_k\) 的列换成 \(V\) 的指标为 \(j_k\) 的列（对 \(1 \leqslant k \leqslant p\)）后所得的行列式；类似地，令 \(\Delta_{K,H}\) 为在 \(V\) 中把指标为 \(j_k\) 的列换成 \(U\) 的指标为 \(i_k\) 的列（对 \(1 \leqslant k \leqslant p\)）后所得的行列式。证明：对每个子集 \(H\) （属于 \([1, n]\) ，具有 \(p\) 个元素），
\end{enumerate}

\[
\Gamma \Delta = \sum_ {\mathbf {K}} \Gamma_ {\mathbf {H}, \mathbf {K}} \Delta_ {\mathbf {K}, \mathbf {H}}
\]

其中 K 遍历 \([1, n]\) 的具有 p 个元素的子集（利用第 6 号中的公式 (21) 和 (22)）。

\begin{enumerate}
\def\labelenumi{\arabic{enumi}.}
\setcounter{enumi}{10}
\tightlist
\item
  设 \(\Delta\) 为方阵 \(X\) （\(n\) 阶）的行列式， \(\Delta_p\) 为方阵 \(\bigwedge^p X\) （其阶为 \(\binom{n}{p}\)）的行列式。证明
\end{enumerate}

\[
\Delta_ {p} \Delta_ {n - p} = \Delta_ {(p)} ^ {(n)}
\]

（利用第 6 号中的公式 (21) 和 (22)）。

\begin{enumerate}
\def\labelenumi{\arabic{enumi}.}
\setcounter{enumi}{11}
\tightlist
\item
  矩阵 \((\alpha_{ij})\) （\(n\) 阶）称为中心对称（相应地，斜中心对称）的，如果 \(\alpha_{n - i + 1, n - j + 1} = \alpha_{ij}\) （相应地，\(\alpha_{n - i + 1, n - j + 1} = -\alpha_{ij}\) ）对 \(1 \leqslant i \leqslant n\) , \(1 \leqslant j \leqslant n\) 成立。
\end{enumerate}

\begin{enumerate}
\def\labelenumi{(\alph{enumi})}
\item
  证明 2p 偶数阶中心对称矩阵的行列式可以表示为两个 p 阶行列式的乘积之形式，而 \(2p + 1\) 阶中心对称矩阵的行列式可以表示为一个 p 阶行列式与一个 \(p + 1\) 阶行列式的乘积之形式。
\item
  证明 \(2p\) 偶数阶斜中心对称矩阵的行列式可以表示为两个 \(p\) 阶行列式的乘积之形式。\(2p + 1\) 奇数阶斜中心对称矩阵的行列式，若在 A 中关系 \(2\xi = 0\) 蕴含 \(\xi = 0\) ，则为零；否则它可以表示为 \(\alpha_{p+1,p+1}\) 与两个 \(p\) 阶行列式的乘积之形式。
\end{enumerate}

\begin{enumerate}
\def\labelenumi{\arabic{enumi}.}
\setcounter{enumi}{12}
\tightlist
\item
  设 \(\Delta = \det(\alpha_{ij})\) 是一个 \(n\) 阶行列式， \(\Delta_{ij}\) 是 \(n - 1\) 阶子式，即 \(\alpha_{ij}\) 的余子式。证明
\end{enumerate}

\[
\left| \begin{array}{c c c c c} \alpha_ {1 1} & \alpha_ {1 2} & \ldots & \alpha_ {1 n} & x _ {1} \\ \alpha_ {2 1} & \alpha_ {2 2} & \ldots & \alpha_ {2 n} & x _ {2} \\ \cdot & \cdot & \cdot & \cdot & \cdot \\ \alpha_ {n 1} & \alpha_ {n 2} & \ldots & \alpha_ {n n} & x _ {n} \\ y _ {1} & y _ {2} & \ldots & y _ {n} & z \end{array} \right| = \Delta z - \sum_ {i, j} (- 1) ^ {i + j} \Delta_ {i j} x _ {i} y _ {j}.
\]

若 \(\Delta = 0\) 且诸 \(\alpha_{ij}\) 属于一个域，证明上述行列式是 \(x_1, x_2, \ldots, x_n\) 的一个线性形式与 \(y_1, y_2, \ldots, y_n\) 的一个线性形式的乘积（利用 §5 习题 5 与 II §6 习题 8）。

给出一个例子，说明当标量环 A 不是域时上述结果不再成立（取 A 为环 \(\mathbf{Z} / (6)\) 且 \(n = 2\)）。

\begin{enumerate}
\def\labelenumi{\arabic{enumi}.}
\setcounter{enumi}{13}
\tightlist
\item
  证明恒等式
\end{enumerate}

\[
\left| \begin{array}{cccccc} 0 & 1 & 1 & 1 & \ldots & 1 \\ 1 & 0 & a _ {1} + a _ {2} & a _ {1} + a _ {3} & \ldots & a _ {1} + a _ {n} \\ 1 & a _ {2} + a _ {1} & 0 & a _ {2} + a _ {3} & \ldots & a _ {2} + a _ {n} \\ 1 & a _ {3} + a _ {1} & a _ {3} + a _ {2} & 0 & \ldots & a _ {3} + a _ {n} \\ \cdot & \cdot & \cdot & \cdot & \cdot & \cdot \\ 1 & a _ {n} + a _ {1} & a _ {n} + a _ {2} & a _ {n} + a _ {3} & \ldots & 0 \\ \end{array} \right| = (- 1) ^ {n} 2 ^ {n - 1} \sum_ {i = 1} ^ {n} a _ {1} a _ {2} \ldots a _ {i - 1} a _ {i + 1} \ldots a _ {n}
\]

（利用习题 13）。

\begin{enumerate}
\def\labelenumi{\arabic{enumi}.}
\setcounter{enumi}{14}
\item
  证明：若 A 上的 n 阶方阵的列线性无关，则该矩阵的行也线性无关（参见 II，§10，习题 3）。
\item
  设 E, F 为两个自由 A-模，维数分别为 m 和 n。为使线性映射 \(u: E \rightarrow F\) 是单射，必须且只需 \(m \leqslant n\) ，并且，若 X 表示 u 关于 E 和 F 的任意两组基的矩阵，则不存在标量 \(\mu \neq 0\) ，使得 X 的所有 m 阶子式乘以 \(\mu\) 后均为零。
\end{enumerate}

\begin{enumerate}
\def\labelenumi{\arabic{enumi}.}
\setcounter{enumi}{16}
\tightlist
\item
  设
\end{enumerate}

\[
\sum_ {j = 1} ^ {n} \alpha_ {i j} \xi_ {j} = \beta_ {i} \quad (1 \leqslant i \leqslant m)\tag{*}
\]

是交换环 A 上含 n 个未知数的 m 个线性方程组成的方程组。令 \(x_{j}\) ( \(1 \leqslant j \leqslant n\) ）为该方程组的矩阵 \(X = (\alpha_{ij})\) 的列，而 \(y = (\beta_{i})\) ；假设在 X 中所有阶数 \textgreater p 的子式均为零，但 \(x_{1} \wedge x_{2} \wedge \cdots \wedge x_{p} \neq 0\) 。为使方程组 (*) 有解，必须 \(x_{1} \wedge x_{2} \wedge \cdots \wedge x_{p} \wedge y = 0\) 。反之，若此条件成立，则存在 \(n + 1\) 个属于 A 的元素 \(\xi_{j}\) ( \(1 \leqslant j \leqslant n + 1\) )，使得 \(\xi_{n+1} \neq 0\) 且

\[
\sum_ {j = 1} ^ {n} \alpha_ {i j} \xi_ {j} = \beta_ {i} \xi_ {n + 1} \quad (1 \leqslant i \leqslant m).
\]

\begin{enumerate}
\def\labelenumi{\arabic{enumi}.}
\setcounter{enumi}{17}
\tightlist
\item
  设 E 和 F 分别为维数 m 和 n 的两个自由 A-模， \(u:E \rightarrow F\) 是一个线性映射，而 X 是 u 关于 E 和 F 的任意两组基的矩阵。
\end{enumerate}

\begin{enumerate}
\def\labelenumi{(\alph{enumi})}
\item
  若 \(m \geqslant n\) 并且存在 \(X\) 的一个 \(n\) 阶可逆子式，则 \(u\) 是满射。
\item
  反之，若 \(u\) 是满射，证明 \(m \geqslant n\) ，且存在 \(X\) 的一个非零的 \(n\) 阶子式。此外，若在环 \(A\) 中，由不可逆元素生成的理想异于 \(A\) ，则存在 \(X\) 的一个 \(n\) 阶可逆子式（考虑外幂 \(\bigwedge^{n}u\) ）。
\item
  设 B 是一个交换环，A 是乘积环 \(B \times B\) 。给出一个从 A-模 \(A^{2}\) 到 A 上的满射线性映射的例子，其矩阵（关于 \(A^{2}\) 与 A 的典范基）的所有元素都是零因子。
\end{enumerate}

\begin{enumerate}
\def\labelenumi{\arabic{enumi}.}
\setcounter{enumi}{18}
\item
  设 X 是交换域上的一个矩阵。要使 X 的秩为 p，只需存在 X 的一个 p 阶子式，它 \(\neq0\) ，并且所有包含这个 p 阶子式的 \(p+1\) 阶子式均为零（证明此时 X 的每一列都是所述 p 阶子式的元素所属的那 p 列的线性组合）。
\item
  设 \(\mathbf{A}\) 为一个交换环， \(\mathbf{M}\) 为任意一个 A-模，而 \(U = (\alpha_{ij})\) 为一个类型为 \((m, n)\) 、元素属于 \(\mathbf{A}\) 的矩阵。
\end{enumerate}

\begin{enumerate}
\def\labelenumi{(\alph{enumi})}
\item
  假设 \(m = n\) ，并设 \(\Delta = \det(U)\) ；那么，若 \(x_1, \ldots, x_n\) 是 \(M\) 的元素，使得 \(\sum_{j=1}^{n} \alpha_{ij} x_j = 0\) （对 \(1 \leqslant i \leqslant n\) ），则 \(\Delta x_i = 0\) 对每个指标 \(i\) 成立。
\item
  假设 \(m\) 和 \(n\) 是任意的， \(U\) 的所有阶 \(> r\) 的子式均为零，且存在一个 \(r\) 阶子式是可逆的，例如矩阵 \((\alpha_{ij})\) （其中 \(i \leqslant r\) 且 \(j \leqslant r\) ）的子式；那么，若 \(a_i\) 表示指标为 \(i\) 的行（在 \(U\) 中），则这些 \(a_i\) 中指标 \(\leqslant r\) 的那些构成 \(A^n\) 的由 \(U\) 的所有行生成的子模的一组基；因而对于 \(i > r\) ，有 \(a_i = \sum_{k=1}^{r} \lambda_{ik} a_k\) 。然后设 \(y_1, \ldots, y_m\) 是 \(M\) 的元素；为使存在元素 \(x_1, \ldots, x_n\) （属于 \(M\) ）满足方程组 \(\sum_{j=1}^{n} \alpha_{ij} x_j = y_i\) （对 \(1 \leqslant i \leqslant m\) ），必要且充分的条件是 \(y_i = \sum_{k=1}^{r} \lambda_{ik} y_k\) （对 \(i > r\) ）。
\end{enumerate}

¶ 21. 设 K 为一个无限交换域，m、n、r 为三个整数 \(\geqslant0\) ，使得 \(r\leqslant n\leqslant m\) ；设 \(\mathbf{M}_{m,n}(\mathbf{K})\) 为由 K 上类型为 \((m,n)\) 的矩阵组成的（mn 维）K-向量空间，而 V 为 \(\mathbf{M}_{m,n}(\mathbf{K})\) 的一个向量子空间，使得 r 是矩阵 \(X\in V\) 的秩的最大值。我们打算证明不等式

\[
\dim (\mathbf {V}) \leqslant m r.\tag{*}
\]

\begin{enumerate}
\def\labelenumi{(\alph{enumi})}
\tightlist
\item
  证明我们可以假设 \(m = n\) 且矩阵
\end{enumerate}

\[
X _ {0} = \left( \begin{array}{c c} I _ {r} & 0 \\ 0 & 0 \end{array} \right)
\]

属于 V。由此推出每个矩阵 \(X \in \mathbf{V}\) 都具有形式

\[
\left( \begin{array}{c c} X _ {1 1} & X _ {1 2} \\ X _ {2 1} & 0 \end{array} \right)
\]

（其中 \(X_{11}\) 的阶为 \(r\) ）且满足条件 \(X_{21}X_{12} = 0\) （利用如下事实：对所有 \(\xi \in \mathbf{K}\) ，矩阵 \(\xi X_0 + X\) 的秩 \(\leqslant r\) ）。

\begin{enumerate}
\def\labelenumi{(\alph{enumi})}
\setcounter{enumi}{1}
\tightlist
\item
  从 (a) 推出，对于两个矩阵 \(X\) , \(Y\) （属于 \(\mathbf{V}\) ），
\end{enumerate}

\[
X _ {2 1} Y _ {1 2} + Y _ {2 1} X _ {1 2} = 0.
\]

\begin{enumerate}
\def\labelenumi{(\alph{enumi})}
\setcounter{enumi}{2}
\tightlist
\item
  对每个矩阵 \(X \in \mathbf{V}\) ，我们记 \(F(X) = (X_{11}X_{12}) \in \mathbf{M}_{r,n}(\mathbf{K})\) ；\(f\) 是一个线性映射，且 \(\dim(\mathbf{V}) = \dim(\operatorname{Im}(f)) + \dim(\operatorname{Ker}(f))\) 。另一方面，设 \(u_x\) 表示线性形式 \((\mathrm{Y}_{11}\mathrm{Y}_{12}) \mapsto \operatorname{Tr}(\mathrm{X}_{21}\mathrm{Y}_{12})\) （定义在 \(\mathbf{M}_{r,n}(\mathbf{K})\) 上）。证明线性映射 \(X \mapsto u_X\) （从 \(\operatorname{Ker}(f)\) 到对偶 \((\mathbf{M}_{r,n}(\mathbf{K}))^*\) 中）是单射；并通过注意在 \(X \mapsto u_X\) 下 \(\operatorname{Ker}(f)\) 的像包含于 \(\operatorname{Im}(f)\) 的正交补中来完成 (*) 的证明。
\end{enumerate}

¶ 22. 设 F 是从 \(\mathbf{M}_{n}(\mathrm{K})\) 到一个集合 E 中的映射，其中 K 是一个交换域，使得对每三个矩阵 \(X\) , \(Y\) , \(Z\) （属于 \(\mathbf{M}_{n}(\mathrm{K})\) ），

\[
\mathbf {F} (X Y Z) = \mathbf {F} (X Z Y).
\]

我们排除 \(\mathbf{M}_{2}(\mathbf{F}_{2})\) 的情形；然后证明存在一个从 K 到 E 中的映射 \(\Phi\) ，使得 \(\mathrm{F}(X)=\Phi(\det(X))\) 。（利用命题 17（第 9 号）的推论，首先证明，对于 \(U\in\mathbf{SL}_{n}(\mathrm{K})\) ，有 \(\mathrm{F}(X)=\mathrm{F}(XU)\) ；由此推出，在 \(\mathbf{GL}_{n}(\mathrm{K})\) 中， \(\mathrm{F}(X)\) 仅依赖于 \(\det(X)\) 。另一方面，若对于 \(r\leqslant n-1\) ，我们记

\[
X _ {r} = \left( \begin{array}{c c} I _ {r} & 0 \\ 0 & 0 \end{array} \right)
\]

，证明存在一个矩阵 \(Y\) （秩为 \(r\) ），使得 \(X_{r-1} = YX_r\) 且 \(Y = X_rY\) ，并由此推出 \(F(X)\) 对所有秩 \(< n\) 的矩阵取相同的值。

\begin{enumerate}
\def\labelenumi{\arabic{enumi}.}
\setcounter{enumi}{22}
\tightlist
\item
  在每个A-代数E中，若 \(x_{1}, \ldots, x_{n}\) 是E的元素，我们记作
\end{enumerate}

\[
\left[ x _ {1} x _ {2} \dots x _ {n} \right] = \sum_ {\sigma \in \mathfrak {S} _ {n}} \varepsilon_ {\sigma} x _ {\sigma (1)} x _ {\sigma (2)} \dots x _ {\sigma (n)}.
\]

证明：若 E 在 A 上具有有限生成系，则存在整数 N，使得 \([x_{1}x_{2}\ldots x_{N}] = 0\) 对 E 的所有元素 \(x_{j} (1 \leqslant j \leqslant N)\) 成立。

\begin{enumerate}
\def\labelenumi{\arabic{enumi}.}
\setcounter{enumi}{23}
\tightlist
\item
  设 \(X_{i}\) ( \(1 \leqslant i \leqslant m\) ) 是交换环 A 上具有相同阶 \(n\) 的方阵。证明若 \(m\) 是偶数，则（习题 23）
\end{enumerate}

\[
\operatorname{Tr} \left[ X _ {1} X _ {2} \dots X _ {m} \right] = 0
\]

且若 \(m\) 是奇数，则

\[
\operatorname{Tr} \left[ X _ {1} X _ {2} \dots X _ {m} \right] = m. \operatorname{Tr} \left(X _ {m} \left[ X _ {1} X _ {2} \dots X _ {m - 1} \right]\right).
\]

（若 \(\lambda\) 是循环置换 \((1,2,\ldots,m)\) ，C 是 \(\mathfrak{S}_m\) 的由 \(\lambda\) 生成的循环子群，而 H 是 \(\mathfrak{S}_m\) 的由保持 \(m\) 不变的置换组成的子群，注意 \(\mathfrak{S}_m\) 中的每个置换都可以唯一地写成 \(\sigma\tau\) ，其中 \(\sigma\in\mathrm{H}\) 且 \(\tau\in\mathrm{C}\) ；然后利用矩阵乘积的迹在因子的循环置换下不变这一事实。）

¶ 25. 设 A 是一个包含有理数域 Q 的交换环，并设 X 是 A 上一个偶数阶 \(2n \geqslant 2\) 的方阵。设 L 为子集 \(\{1, 2, \ldots, n + 1\}\) （属于 \(\{1, 2, \ldots, 2n\}\) ），而 L' 为其补集。证明以下公式（“Redei 恒等式”）：

\[
(*) \quad \det (X _ {\mathrm{L}, \mathrm{L}}) \det (X _ {\mathrm{L} ^ {\prime}, \mathrm{L} ^ {\prime}}) = \sum_ {\mathrm{H}, \mathrm{K}} (- 1) \frac {r - 1}{n \binom {n - 1} {r}} \det (X _ {\mathrm{H}, \mathrm{K}}) \det (X _ {\mathrm{H} ^ {\prime}, \mathrm{K} ^ {\prime}})
\]

其中 H 遍历 \(\{1,2,\ldots,2n\}\) 的包含 l 的 n 元子集组成的集合，而 K 遍历 L 的 n 元子集组成的集合，且 \(r=\operatorname{Card}(H\cap L')\) 。（计算 (*) 右端中项 \(x_{\sigma(1),1}x_{\sigma(2),2}\ldots x_{\sigma(2n),2n}\) 对任意置换 \(\sigma\in S_{2n}\) 的系数，注意该系数仅可能 \(\neq0\) （当 \(H=\sigma(K)\) ）；证明它仅可能 \(\neq0\) （当 \(\sigma(L)=L\) ）。）

\begin{enumerate}
\def\labelenumi{\arabic{enumi}.}
\setcounter{enumi}{25}
\tightlist
\item
  \begin{enumerate}
  \def\labelenumii{(\alph{enumii})}
  \tightlist
  \item
    在命题 19（第 10 号）的记号下，假设存在两个 A{[}X{]}-线性映射 \(g_1, g_2\) （从 M{[}X{]} 到 M'{[}X{]} 中），使得 \(\psi' \circ g_2 = g_1 \circ \psi\) ；证明此时存在一个 A{[}X{]}-线性映射 \(g\) （从 \(M_u\) 到 \(M_{u'}\) 中），使得 \(\phi' \circ g_1 = \phi' \circ \bar{g}\) 。进一步，若 \(g_1\) 和 \(g_2\) 是 A{[}X{]}-同构（从 M{[}X{]} 到 M'{[}X{]} 上），则 \(g\) 是一个 A{[}X{]}-同构（从 \(M_u\) 到 \(M_{u'}\) 上）。
  \end{enumerate}
\end{enumerate}

\begin{enumerate}
\def\labelenumi{(\alph{enumi})}
\setcounter{enumi}{1}
\tightlist
\item
  在第 10 号的记号下，自同态 \(u, u'\) 称为等价的，如果存在两个从 M 到 M 上的同构 \(f_1, f_2\) ，使得 \(u' \circ f_2 = f_1 \circ u\) ；若存在一个从 M 到 M' 上的同构 \(f\) ，使得 \(u' \circ f = f \circ u\) ，则称它们为相似的。由 (a) 推出，要使 \(u\) 和 \(u'\) 相似，当且仅当自同态 \(X - \bar{u}\) 和 \(X - \bar{u}'\) （分别为 A{[}X{]}-模 M{[}X{]} 和 M'{[}X{]} 的自同态）等价。
\end{enumerate}

